\scanpage{151}
\uline{Corollaire 1} Fixons un $\delta=\{\delta_0>\delta_1\cdots>\delta_s>0\}$,
considérons $D_{\delta_0,\delta_1,\ldots,\delta_s}=D_\delta$, le diviseur
$\Theta_\delta$ sur $D_\delta$, et $V^*_{\Theta_\delta,D_\delta}=D_\delta^{*\infty}$
(structure à l'infini de $D_\delta^{*\infty}$).

Alors le syst. $\frac12$ simplicial des
\[
\deleted{\unsure{D}}\;V^*_{\underbrace{D_{\underbrace{d_0,\ldots,d_r},\delta_0,\delta_1,\ldots,\delta_s}}_{\text{de degré simp. }r}}
\qquad(\text{pour }d_0>\cdots>d_r>\delta_0
\]
variables) a comme \uline{lim} toposique
la ``structure à l'$\infty$'' $V^*_{\Theta_\delta,D_\delta}=D_\delta^*$.

On applique le th. à $D_\delta,\Theta_\delta\ldots$

\deleted{\uline{Corollaire 2} Supposons que
\diagram{\begin{tikzcd}[ampersand replacement=\&]
V^*_{\Theta,X} \arrow[r,leftrightarrow] \arrow[d,equal,"\mathrm{def}"'] \& X^* \arrow[d,equal] \\
X^{*\infty} \& X\setminus\Theta
\end{tikzcd}}
soit un $1$-\unsure{équivalence} (i.e induise un iso. pour les $\pi_0$, et les $\pi_1$ des comp. conn. correspondantes) alors on a
\[
\deleted{\unsure{\widetilde\Pi_{X,X^*}}}\simeq
\underset{\longrightarrow}{\operatorname{Lim}}\bigl(\deleted{\unsure{\widetilde\Pi_{d_0},\widetilde\Pi}}\bigr)
\quad
\begin{array}{ll}
\widetilde\Pi_{d_0}\deleted{\unsure{\Pi_{d_0}}}&(1\leq d_0\leq3)\\
\widetilde\Pi_{d_0,d_1}&(d_0>d_1\geq
\end{array}
\]
}
\margin{150}
